\documentclass[12pt]{article}
\usepackage[utf8]{inputenc}
\usepackage[T1]{fontenc}
\usepackage[margin=1in]{geometry}
\usepackage{setspace}
\usepackage{xcolor}
\usepackage[colorlinks=true, urlcolor=blue!60!black, linkcolor=blue!60!black]{hyperref}
\setlength{\parindent}{2em}
\setlength{\parskip}{1.25em}
\renewcommand{\baselinestretch}{1.0}
\begin{document}
\begin{center}
{\LARGE \textit{Ashcroft v.\ Free Speech Coalition}, 535 U.S.\ 234 (2002)\par}
\vspace{0.8em}
{\large Case Brief\par}
\end{center}
\vspace{1em}

\section{Memory Jogger}

In \href{https://supreme.justia.com/cases/federal/us/535/234/}{\textit{Ashcroft v.\ Free Speech Coalition}}, the Supreme Court struck down the Child Pornography Prevention Act of 1996's bans on sexually explicit images that merely ``appear to be'' of minors or ``convey the impression'' of depicting minors as substantially overbroad under the First Amendment, holding that government may not suppress non-obscene speech created without harming real children.

\section{Detailed Case Facts}

Before 1996, federal law (18 U.S.C.\ \S~2252) defined child pornography in line with \href{https://supreme.justia.com/cases/federal/us/458/747/}{\textit{New York v.\ Ferber}}, 458 U.S.\ 747 (1982), covering only sexually explicit visual depictions produced using actual minors. Concerned by rapid advances in computer-imaging technology during the 1990s, Congress enacted the Child Pornography Prevention Act of 1996 (CPPA), Pub.\ L.\ No.\ 104--208, \S~121, 110 Stat.\ 3009--26, which retained the prohibition on actual child pornography in 18 U.S.C.\ \S~2256(8)(A) and added three new categories, imposing severe felony penalties of up to 15 years' imprisonment for a first non-possession offense (and 5 to 30 years for repeat offenders) and up to 5 years for possession (18 U.S.C.\ \S~2252A(b)).

Two of those new provisions were challenged in this litigation:

\begin{itemize}
\item \textbf{The ``Appears to Be'' Provision (18 U.S.C.\ \S~2256(8)(B)):} Prohibited ``any visual depiction, including any photograph, film, video, picture, or computer or computer-generated image or picture,'' that ``is, or appears to be, of a minor engaging in sexually explicit conduct.'' This swept in both wholly computer-generated ``virtual child pornography'' and ``youthful-looking adult'' pornography, regardless of how the image was made and without requiring the work to be obscene under \textit{Miller v.\ California}, 413 U.S.\ 15 (1973).
\item \textbf{The ``Conveys the Impression'' (Pandering) Provision (18 U.S.C.\ \S~2256(8)(D)):} Prohibited any sexually explicit visual depiction that is ``advertised, promoted, presented, described, or distributed in such a manner that conveys the impression'' that it depicts ``a minor engaging in sexually explicit conduct.'' Once pandered in that way, the material remained contraband even in the hands of subsequent possessors who took no part in the marketing.
\end{itemize}

A third new provision, \S~2256(8)(C), prohibited computer ``morphing''---altering innocent photographs of real children to make them appear engaged in sexual conduct. Because morphed images implicate the interests of real children, the respondents did not challenge \S~2256(8)(C), and the Court did not pass upon it. The CPPA also contained an affirmative defense (\S~2252A(c)) allowing defendants charged with non-possession offenses to avoid conviction by proving that the images were produced using only actual adults and were not pandered as child pornography; however, that defense did not apply to simple possession or to computer-generated images.

In congressional findings accompanying the CPPA (notes following 18 U.S.C.\ \S~2251), Congress justified banning virtual depictions on four grounds: (1) pedophiles might use virtual images to seduce real children; (2) such images ``whet their own sexual appetites,'' leading to abuse of actual children; (3) virtual images fuel the market for real child pornography; and (4) increasingly realistic computer graphics make it difficult for prosecutors to prove beyond a reasonable doubt that an image depicts a real child.

Fearing felony prosecution, the Free Speech Coalition (a California adult-entertainment trade association whose members did not use minors), Bold Type, Inc.\ (publisher of a book advocating the nudist lifestyle), Jim Gingerich (a painter of nudes), and Ron Raffaelli (an erotic photographer) filed a pre-enforcement facial challenge claiming that \S\S~2256(8)(B) and 2256(8)(D) were unconstitutionally overbroad and vague under the First Amendment.

\section{Procedural History}

In the U.S.\ District Court for the Northern District of California, Judge Samuel Conti rejected the plaintiffs' First Amendment claims and granted summary judgment to the Government (\textit{The Free Speech Coalition v.\ Reno}, No.\ C 97--0281 VSC, 1997 WL 487758 (N.D.\ Cal.\ Aug.\ 12, 1997)), reasoning that it was ``highly unlikely'' that legitimate adaptations of works like \textit{Romeo and Juliet} would ever be prosecuted as contraband.

The U.S.\ Court of Appeals for the Ninth Circuit reversed (\textit{The Free Speech Coalition v.\ Reno}, 198 F.\ 3d 1083 (9th Cir.\ 1999)), holding that \S\S~2256(8)(B) and 2256(8)(D) were facially overbroad and void for vagueness because they banned speech that was neither obscene under \textit{Miller} nor produced through the exploitation of real children under \textit{Ferber}. Judge Warren J.\ Ferguson dissented, arguing that virtual child pornography should be recognized as a new category of unprotected speech. The Ninth Circuit denied rehearing en banc over the dissent of three judges (\href{https://law.justia.com/cases/federal/appellate-courts/F3/220/1113/625466/}{220 F.\ 3d 1113 (9th Cir.\ 2000)}).

Because four other circuits---the First, Fourth, Fifth, and Eleventh Circuits---had upheld the CPPA against facial First Amendment challenges, the Supreme Court granted certiorari on January 22, 2001 (531 U.S.\ 1124). Oral argument was held on October 30, 2001 (with Deputy Solicitor General Paul D.\ Clement arguing for Attorney General John Ashcroft and H.\ Louis Sirkin arguing for respondents), and the Court issued its decision on April 16, 2002, affirming the Ninth Circuit (\href{https://www.oyez.org/cases/2001/00-795}{Oyez}; \href{https://www.law.cornell.edu/supct/html/00-795.ZS.html}{Cornell LII}).

\section{Judicial Votes}

The Court affirmed the Ninth Circuit's judgment invalidating both provisions by a \textbf{6--3} vote on the core virtual-pornography ban (\S~2256(8)(B)) and \textbf{7--2} on both youthful-adult pornography and the pandering provision (\S~2256(8)(D)):

\begin{itemize}
\item \textbf{Majority Opinion (5 Justices):} Justice Anthony M.\ Kennedy delivered the opinion of the Court, joined by Justices John Paul Stevens, David H.\ Souter, Ruth Bader Ginsburg, and Stephen G.\ Breyer. (Struck down both \S~2256(8)(B) and \S~2256(8)(D) in full as substantially overbroad.)
\item \textbf{Concurring in the Judgment (Justice Thomas):} Justice Clarence Thomas concurred in the judgment, agreeing that the CPPA failed on the current record while leaving open narrower future regulation if technology actually thwarts prosecutions of real child pornography.
\item \textbf{Concurring in the Judgment in Part and Dissenting in Part (Justice O'Connor):} Justice Sandra Day O'Connor concurred in striking down \S~2256(8)(D) and \S~2256(8)(B) as applied to youthful-looking adult pornography, but dissented as to computer-generated virtual child pornography under \S~2256(8)(B). Chief Justice William H.\ Rehnquist and Justice Antonin Scalia joined Part II of her opinion.
\item \textbf{Dissenting Opinion (2 Justices):} Chief Justice William H.\ Rehnquist filed a dissenting opinion, joined by Justice Antonin Scalia except as to the single paragraph discussing legislative history. (Would have upheld both provisions through narrowing constructions.)
\end{itemize}

\section{Holding}

\textbf{Judgment:} Affirmed (holding 18 U.S.C.\ \S\S~2256(8)(B) and 2256(8)(D) facially overbroad and unconstitutional under the First Amendment).

\textbf{Rule Announced:} Speech that is neither obscene under \textit{Miller v.\ California}, 413 U.S.\ 15 (1973), nor produced through the sexual exploitation of actual children under \href{https://supreme.justia.com/cases/federal/us/458/747/}{\textit{New York v.\ Ferber}}, 458 U.S.\ 747 (1982), retains First Amendment protection. Congress may not categorically ban lawful, non-obscene visual depictions of fictional or adult-portrayed teenage sexuality on the theory that they might indirectly encourage child abuse or make prosecuting real child pornography harder; doing so suppresses a substantial amount of protected expression in violation of the First Amendment overbreadth doctrine (\href{https://supreme.justia.com/cases/federal/us/413/601/}{\textit{Broadrick v.\ Oklahoma}}, 413 U.S.\ 601 (1973)). Likewise, a pandering statute is overbroad if it criminalizes mere possession of non-obscene adult works simply because someone earlier in the distribution chain marketed them as child pornography.

\textbf{Holding vs.\ Dicta:} The facial invalidation of \S\S~2256(8)(B) and 2256(8)(D) as substantially overbroad---and the Court's refusal to reach the respondents' void-for-vagueness claim---constitute the holding. The Court's observations that morphed images of real children (\S~2256(8)(C)) ``implicate the interests of real children and are in that sense closer to the images in \textit{Ferber},'' that an affirmative defense might or might not save a narrower statute, and that virtual images might actually displace rather than stimulate the market for real child pornography are non-binding dicta that shaped subsequent federal legislation, including the PROTECT Act of 2003.

\section{Analysis of Opinions}

\subsection{Majority Opinion (Justice Kennedy)}

Writing for the Court, Justice Kennedy began with the bedrock First Amendment principle that severe criminal penalties on speech chill protected expression, making facial overbreadth review essential under \href{https://supreme.justia.com/cases/federal/us/413/601/}{\textit{Broadrick v.\ Oklahoma}}, 413 U.S.\ 601, 612 (1973). Because the CPPA targeted speech outside the narrow historic categories of unprotected expression, the Government had to justify its bans under either \textit{Miller}, \textit{Ferber}, or an independent compelling interest. Kennedy systematically dismantled each theory:

\begin{itemize}
\item \textbf{Inconsistency with \textit{Miller}:} Under \textit{Miller v.\ California}, 413 U.S.\ at 24, obscenity requires proof that a work, taken as a whole, appeals to the prurient interest, is patently offensive under community standards, and lacks serious literary, artistic, political, or scientific value. The CPPA dispensed with all three prongs. It banned any single sexually explicit scene depicting someone who appears under 18---an age higher than the legal age of consent in most States---thereby sweeping in Renaissance paintings, modern film adaptations of Shakespeare's \textit{Romeo and Juliet}, and acclaimed movies like \textit{Traffic} and \textit{American Beauty}, violating the rule that a work's artistic merit must be judged as a whole.
\item \textbf{No Support in \textit{Ferber} or \textit{Osborne}:} \href{https://supreme.justia.com/cases/federal/us/458/747/}{\textit{Ferber}} and \href{https://supreme.justia.com/cases/federal/us/495/103/}{\textit{Osborne v.\ Ohio}}, 495 U.S.\ 103 (1990), permitted bans on non-obscene child pornography only because the speech was ``intrinsically related'' to the sexual abuse of real children---serving as a permanent record of a child's abuse and an economic motive for its production. By contrast, virtual child pornography ``records no crime and creates no victims by its production.'' Indeed, \textit{Ferber} itself relied on youthful-looking adults and simulation as lawful, constitutionally protected alternatives for artistic works (458 U.S.\ at 763); the CPPA could not rely on \textit{Ferber} to outlaw the very alternatives \textit{Ferber} invoked.
\item \textbf{Rejection of Indirect-Harm and Enforcement Rationales:} First, the claim that pedophiles might use virtual images to seduce children violates \textit{Butler v.\ Michigan}, 352 U.S.\ 380 (1957), and \textit{Sable Communications of California, Inc.\ v.\ FCC}, 492 U.S.\ 115 (1989): government may not reduce adults to viewing only what is fit for children, and it may directly punish grooming or solicitation instead. Second, the claim that virtual images ``whet the appetites'' of pedophiles clashes with \textit{Stanley v.\ Georgia}, 394 U.S.\ 557, 566 (1969), and \textit{Brandenburg v.\ Ohio}, 395 U.S.\ 444, 447 (1969): the mere tendency of speech to encourage unlawful thoughts or indefinite future conduct cannot justify suppression absent a direct link to imminent lawless action. Third, the market-deterrence theory was both economically implausible (indistinguishable virtual substitutes would undercut rather than boost demand for risky real abuse) and legally insufficient where no underlying crime occurred. Fourth, banning protected speech because it resembles unprotected speech ``turns the First Amendment upside down.'' Nor did \S~2252A(c)'s affirmative defense save the statute: shifting the burden of proof to speakers raises grave First Amendment concerns, and the defense was unavailable for mere possession or for computer-generated images.
\item \textbf{Invalidation of \S~2256(8)(D) (Pandering):} Unlike \textit{Ginzburg v.\ United States}, 383 U.S.\ 463 (1966), which treated commercial pandering as evidentiary context in close obscenity cases, \S~2256(8)(D) required no commercial exploitation and permanently tainted a sexually explicit film in the hands of downstream possessors who played no role in mislabeling it. It was therefore substantially overbroad.
\end{itemize}

\subsection{Concurring in the Judgment (Justice Thomas)}

Justice Thomas agreed that the CPPA was unconstitutional on the existing record, but distanced himself from the majority's categorical statement that government may never restrict virtual imagery to enforce laws against real child pornography. Identifying the ``prosecution rationale''---that defendants might defeat real-child-pornography prosecutions by claiming images are computer generated---as the Government's strongest argument, Thomas noted that the Government had cited not a single acquittal resulting from such a defense. Yet if future technology truly rendered real child pornography unprosecutable, Thomas would leave Congress room to enact a narrowly tailored virtual-image restriction backed by a complete affirmative defense or other safeguard.

\subsection{Concurring in the Judgment in Part and Dissenting in Part (Justice O'Connor)}

Justice O'Connor agreed with the majority (in Part I) that virtual and youthful-adult pornography are not categorically unprotected under \textit{Ferber}, that \S~2256(8)(D)'s pandering ban fails strict scrutiny, and that \S~2256(8)(B) is overbroad as applied to youthful-looking adult pornography (threatening films like \textit{Traffic} and \textit{American Beauty}).

In Part II (joined by Chief Justice Rehnquist and Justice Scalia), however, O'Connor argued that \S~2256(8)(B)'s ban on computer-generated virtual child pornography satisfies strict scrutiny and is neither vague nor overbroad. Citing realistic computer-generated films such as \textit{Final Fantasy: The Spirits Within} (2001) and \textit{Turner Broadcasting System, Inc.\ v.\ FCC}, 520 U.S.\ 180 (1997), she urged deferring to Congress's predictive judgment about prosecutorial evasion. Invoking constitutional avoidance (\textit{Crowell v.\ Benson}, 285 U.S.\ 22 (1932)), she would read ``appears to be'' narrowly to mean ``virtually indistinguishable from'' actual children, excluding cartoons and sculptures. Because respondents identified no wholly computer-generated films of serious value swept in by that narrowed ban, she argued in Part III that the Court should sever and strike \S~2256(8)(B) only as applied to youthful-adult pornography, following \textit{United States v.\ Grace}, 461 U.S.\ 171 (1983).

\subsection{Dissenting Opinion (Chief Justice Rehnquist)}

Chief Justice Rehnquist, joined by Justice Scalia (except for one paragraph invoking legislative history), would have upheld the entire CPPA through limiting constructions under \textit{Broadrick}. Joining Part II of Justice O'Connor's opinion on the compelling interest in banning virtually indistinguishable computer-generated images, Rehnquist argued that \S~2256(2)'s definition of ``sexually explicit conduct''---even with the word ``simulated''---reaches only ``hard core'' pornography, not suggestive mainstream scenes like actors ``squirming under a blanket'' in \textit{Romeo and Juliet}, \textit{Traffic}, or \textit{American Beauty} (noting that the latter two films were made and won Oscars while the CPPA was on the books). As to \S~2256(8)(D), Rehnquist would narrow ``conveys the impression'' to reach only the ``sordid business of pandering'' under \textit{Ginzburg}, and would read a strong \textit{mens rea} requirement into possession offenses under \href{https://supreme.justia.com/cases/federal/us/513/64/case.pdf}{\textit{United States v.\ X-Citement Video, Inc.}}, 513 U.S.\ 64 (1994), so innocent possessors of mislabeled films would not be prosecuted.

\section{Examples: Future Applications}

\begin{enumerate}
\item \textbf{Same-Side (Unconstitutional): State Ban on Non-Obscene Anime or CGI Depicting Fictional Teenagers.} A State enacts a felony statute banning the possession or distribution of any computer-animated or hand-drawn depiction of a fictional character who ``appears to be under 18'' engaging in sexually explicit conduct, without incorporating \textit{Miller}'s obscenity prongs. \textit{Analysis:} Squarely unconstitutional under \textit{Free Speech Coalition}. Because no actual child is abused in producing purely fictional animation, the speech falls outside \textit{Ferber} and cannot be banned unless it satisfies all three prongs of \textit{Miller} obscenity.
\item \textbf{Same-Side (Unconstitutional): Federal Ban on Videos Depicting Animal Cruelty.} Congress passes a criminal statute (18 U.S.C.\ \S~48) prohibiting the commercial creation, sale, or possession of depictions of animal cruelty, sweeping in hunting videos alongside ``crush videos.'' \textit{Analysis:} Relying directly on \textit{Free Speech Coalition}'s refusal to create new unprotected speech categories through ad hoc cost-benefit balancing and its robust overbreadth framework, the Supreme Court invalidated \S~48 in \href{https://supreme.justia.com/cases/federal/us/559/460/}{\textit{United States v.\ Stevens}}, 559 U.S.\ 460 (2010).
\item \textbf{Opposite-Side (Constitutional): Narrowly Tailored Ban on Offers to Provide or Obtain Actual Child Pornography.} Congress responds to \textit{Free Speech Coalition} by enacting the PROTECT Act's revised pandering provision (18 U.S.C.\ \S~2252A(a)(3)(B)), criminalizing knowingly advertising, promoting, or soliciting material in a manner that reflects the belief, or is intended to cause another to believe, that the material depicts an \textit{actual} minor engaging in sexually explicit conduct. \textit{Analysis:} Unlike \S~2256(8)(D), this provision punishes only the offer to engage in an illegal transaction (unprotected speech integral to criminal conduct) rather than tainting downstream possession of protected works---and the Supreme Court upheld it 7--2 in \href{https://supreme.justia.com/cases/federal/us/553/285/}{\textit{United States v.\ Williams}}, 553 U.S.\ 285 (2008).
\item \textbf{Opposite-Side (Constitutional): Prosecution of Obscene Virtual Child Pornography Under \textit{Miller}.} Federal prosecutors indict a defendant under 18 U.S.C.\ \S~1466A(a)(1) for distributing photorealistic computer-generated images of young children engaged in explicit sexual acts, where the statute explicitly incorporates \textit{Miller}'s obscenity standard (prurient appeal, patent offensiveness, and lack of serious literary, artistic, political, or scientific value). \textit{Analysis:} Constitutional under \textit{Free Speech Coalition}. Justice Kennedy expressly noted that Congress may proscribe obscene depictions of minors under \textit{Miller} and \href{https://www.law.cornell.edu/constitution-conan/amendment-1/child-pornography}{federal obscenity statutes} even when no real child was used in production.
\item \textbf{Fence-Sitter (Close Call): Generative-AI Deepfakes Trained on Real CSAM Datasets or Photorealistic Synthetic Images.} A defendant uses a generative-AI model trained on datasets containing real child sexual abuse material (CSAM) to produce photorealistic images of wholly nonexistent children, and challenges a federal ban under \textit{Free Speech Coalition}. \textit{Analysis:} This is the defining modern test of \textit{Free Speech Coalition}. On one side, the resulting image depicts no identifiable real child, putting it within Kennedy's ``records no crime and creates no victims'' rule (unless prosecuted as \textit{Miller} obscenity). On the other side, modern generative AI arguably triggers Justice Thomas's concurrence---making real and synthetic CSAM forensically indistinguishable at massive scale---and if the AI model was trained on real CSAM, prosecutors can argue the synthetic output is ``intrinsically related'' to real abuse, resembling morphed imagery under \S~2256(8)(C).
\end{enumerate}

\section{Critique}

\textit{Ashcroft v.\ Free Speech Coalition} stands as a landmark defense of free-speech categoricalism, yet it has attracted sustained critique from both child-protection advocates and methodological critics.

\textbf{Child-Protection, Feminist, and Technological Critiques:} Many progressive, feminist, and child-welfare scholars argued that Justice Kennedy's rigid production-harm test understates the real-world dangers of virtual child pornography. Early commentary in the \href{https://scholarship.law.unc.edu/cgi/viewcontent.cgi?article=4081&context=nclr}{\textit{North Carolina Law Review}} and the \href{https://scholarship.law.cornell.edu/cgi/viewcontent.cgi?article=1126&context=cjlpp}{\textit{Cornell Journal of Law and Public Policy}} contended that the majority dismissed congressional findings on grooming, desensitization, and prosecutorial burdens too cavalierly, and that Kennedy's economic assumption---that virtual images would drive real CSAM out of the market---lacked empirical backing. With the rise of modern generative AI, those critiques have intensified sharply: recent articles in the \href{https://scholarship.law.nd.edu/cgi/viewcontent.cgi?article=1970&context=ndjlepp}{\textit{Notre Dame Journal of Law, Ethics \& Public Policy}}, the \href{https://law.ua.edu/wp-content/uploads/2026/06/48LawPsycholRev169.pdf}{\textit{Law \& Psychology Review}}, and the \href{https://www.americanbar.org/groups/litigation/resources/newsletters/childrens-rights/when-image-is-fake-harm-real/}{American Bar Association's Children's Rights Litigation newsletter} argue that cheap, photorealistic AI-generated CSAM now threatens to overwhelm law enforcement and realize the exact forensic deadlock Justice Thomas warned about in 2002.

\textbf{Originalist, Textualist, and Doctrinal Critiques:} From a conservative and textualist vantage point, critics---mirroring Chief Justice Rehnquist and Justice O'Connor---faulted the majority for wielding facial overbreadth as a blunt instrument rather than applying canonical constitutional avoidance or severability. By invoking hypothetical prosecutions of \textit{Romeo and Juliet} and \textit{American Beauty} that had never occurred during the CPPA's six years on the books, critics argued the Court abandoned \textit{Broadrick}'s requirement that overbreadth be ``substantial'' judged in relation to a statute's plainly legitimate sweep. Conversely, First Amendment theorists such as Joseph Blocher in the \href{https://nyulawreview.org/wp-content/uploads/2018/08/NYULawReview-84-2-Blocher.pdf}{\textit{New York University Law Review}} have highlighted how \textit{Free Speech Coalition} reinforced rigid First Amendment categoricalism (culminating in \textit{Stevens})---a methodology praised by free-speech libertarians and many originalists for constraining judicial balancing, even as it creates persistent friction when legislatures confront novel digital harms unforeseen at the Founding.

\section{Key Quotations}

All quotations below are verified verbatim from the published opinion in \href{https://supreme.justia.com/cases/federal/us/535/234/}{\textit{Ashcroft v.\ Free Speech Coalition}, 535 U.S.\ 234 (2002)}:

\begin{quote}
``In contrast to the speech in \textit{Ferber}, speech that itself is the record of sexual abuse, the CPPA prohibits speech that records no crime and creates no victims by its production. Virtual child pornography is not `intrinsically related' to the sexual abuse of children, as were the materials in \textit{Ferber}.''
\end{quote}
\noindent --- Justice Kennedy, majority opinion, 535 U.S.\ at 250.

\begin{quote}
``First Amendment freedoms are most in danger when the government seeks to control thought or to justify its laws for that impermissible end. The right to think is the beginning of freedom, and speech must be protected from the government because speech is the beginning of thought.''
\end{quote}
\noindent --- Justice Kennedy, majority opinion, 535 U.S.\ at 253.

\begin{quote}
``The Government may not suppress lawful speech as the means to suppress unlawful speech. Protected speech does not become unprotected merely because it resembles the latter. The Constitution requires the reverse.''
\end{quote}
\noindent --- Justice Kennedy, majority opinion, 535 U.S.\ at 255.

\begin{quote}
``While this speculative interest cannot support the broad reach of the CPPA, technology may evolve to the point where it becomes impossible to enforce actual child pornography laws because the Government cannot prove that certain pornographic images are of real children. In the event this occurs, the Government should not be foreclosed from enacting a regulation of virtual child pornography that contains an appropriate affirmative defense or some other narrowly drawn restriction.''
\end{quote}
\noindent --- Justice Thomas, concurring in the judgment, 535 U.S.\ at 259.

\begin{quote}
``The aim of ensuring the enforceability of our Nation's child pornography laws is a compelling one. The CPPA is targeted to this aim by extending the definition of child pornography to reach computer-generated images that are virtually indistinguishable from real children engaged in sexually explicit conduct.''
\end{quote}
\noindent --- Chief Justice Rehnquist, dissenting, 535 U.S.\ at 273.


\end{document}
